% !PRINTED 784
\weberpage{784}
%<<CONT
$s$ pass over into 1, then $\log Q_1(s)$ becomes infinite by 4., while the remaining terms of the left side, $\log Q_2(s), \ldots \log Q_\mu(s)$, remain finite. Since $R(s)$ likewise remains finite, the sum $\Sigma p^{-s}$ must become infinite, and this is certainly possible only if the sum consists of infinitely many terms. With this it is proved:

5. \gesperrt{In each of the number classes $A$ (with respect to the modulus $m$) infinitely many primes are contained.}

One can also express this by saying that the linear form $mx + a$, in which $m$ and $a$ are integers without a common divisor, represents a prime for infinitely many integral values of $x$.

\begin{center}
\rule{4cm}{0.4pt}
\end{center}

% [The remainder of this leaf and the following leaf are blank in the source.]
